\chapter{Cửa Hàng Hỏi Ngu \& Trả Lời Khịa}
\markboth{Cửa Hàng Hỏi Ngu \& Trả Lời Khịa}{Cửa Hàng Hỏi Ngu \& Trả Lời Khịa}

\begin{truyen}{Màn 1–50 trong bộ 1000 màn}{Scenes 1–50 of the 1000-scene collection}
\chuhoa{L}{àm thầy cô không khó, cái khó là giữ bình tĩnh trước những câu hỏi ``đi vào lòng đất'' của học sinh. Chương này ghi lại những màn đối đáp kinh điển --- nơi sự ngây thơ đáng yêu của các em gặp phải sự ``khịa'' duyên dáng của thầy cô. Không ai thắng, không ai thua, chỉ có sự thật hiện ra sau mỗi câu đáp.}
\textit{(Being a teacher isn't hard; the hard part is staying calm when faced with `earth-shattering' questions from students. This chapter records classic exchanges --- where student innocence meets the witty sharp edge of teacher wisdom. No one wins, no one loses; only the truth emerges after each reply.)}
\end{truyen}

%------------------------------------------------------------
\section{Hỏi Về Bài Vở, Điểm Số và Thầy Cô (Màn 1–25)}
%------------------------------------------------------------
\begin{enumerate}[leftmargin=*]

\item
\studentpair{Thầy ơi, sao thầy không cho em điểm cao hơn điểm thực tế?}{Teacher, why don't you give me a grade higher than what I actually scored?}
\teacherpair{Thầy không biết ai dạy em về ``điểm thực tế'', nhưng thầy đảm bảo điểm thực tế của em là cái số thầy vừa viết.}{I'm not sure who taught you about `actual grades,' but I assure you your actual grade is exactly the number I just wrote.}

\item
\studentpair{Em có thể nộp bài trễ không ạ?}{Can I submit the assignment late?}
\teacherpair{Tất nhiên có thể, nhưng điểm thì không thể nộp trễ theo.}{Of course you can --- but your grade cannot be submitted late along with it.}

\item
\studentpair{Câu này khó quá thầy ơi, em không hiểu đề bài.}{This question is too hard --- I don't understand what it's asking.}
\teacherpair{Để thầy diễn đạt lại: em đọc không hiểu, hay em không đọc?}{Let me rephrase the question: you read it and didn't understand, or you didn't read it?}

\item
\studentpair{Thầy có thể cho em coi lại bài của bạn không, để em biết mình làm sai chỗ nào?}{Can I look at my classmate's answers to see where I went wrong?}
\teacherpair{Hoặc em có thể coi lại bài của chính em và tìm ra chỗ sai đó. Thử xem.}{Or you could look at your own answers and find where you went wrong. Give it a try.}

\item
\studentpair{Thầy ơi, tại sao đề thi ra những cái em không học?}{Teacher, why does the exam have things I didn't study?}
\teacherpair{Vì đề thi ra những cái thầy đã dạy, mà thầy dạy hết rồi.}{Because the exam covers what I taught --- and I taught all of it.}

\item
\studentpair{Em học rồi nhưng vào phòng thi quên hết thầy ơi.}{I studied, but I forgot everything in the exam room.}
\teacherpair{Đó là lý do thầy luôn nói học hiểu chứ không phải học thuộc, nhưng em đã không nghe phần đó.}{That is exactly why I always say understand it, don't just memorize it --- but you didn't hear that part.}

\item
\studentpair{Thầy thấy em có tiến bộ gì không ạ?}{Do you see any improvement in me?}
\teacherpair{Có, em đang ngày càng đặt câu hỏi hay hơn về sự tiến bộ của mình.}{Yes --- your questions about your own progress are getting noticeably better.}

\item
\studentpair{Thầy hay gọi em lên bảng có phải vì thầy ghét em không?}{Do you call me to the board because you dislike me?}
\teacherpair{Không phải. Thầy gọi em lên để kiểm tra xem bảng đen có giúp em nhớ được gì không.}{Not at all. I call you up to check whether the blackboard can help you remember anything.}

\item
\studentpair{Thầy ơi, em học ngành gì để không phải thi nhiều?}{What field should I study to avoid lots of exams?}
\teacherpair{Thầy chưa thấy ngành học nào như vậy, nhưng nếu em tìm ra thì nhớ báo thầy với.}{I haven't found one yet --- but if you discover it, please let me know.}

\item
\studentpair{Có cách nào đọc sách mà vẫn nhớ không thầy?}{Is there a way to actually remember what I read?}
\teacherpair{Có, em có thể thử đọc sách thay vì chỉ mở sách ra và nhìn vào trang bìa.}{Yes --- you could try reading the book instead of just opening it and staring at the cover.}

\item
\studentpair{Em không thích môn này thầy ơi.}{I don't like this subject.}
\teacherpair{Thầy hiểu, nhưng môn này không cần em thích để em vẫn phải học và thi.}{I understand --- but this subject doesn't need you to like it in order for you to have to study and be tested on it.}

\item
\studentpair{Em học suốt mà sao vẫn không vô?}{I study all the time but nothing sticks.}
\teacherpair{``Học suốt'' là học liên tục hay cầm sách liên tục? Vì đó là hai thứ rất khác nhau.}{`All the time' --- do you mean continuously studying, or continuously holding the book? Those are two very different things.}

\item
\studentpair{Sao môn này khó vậy thầy?}{Why is this subject so hard?}
\teacherpair{Vì nó đòi hỏi suy nghĩ. Không phải môn nào cũng vậy, nên em mới thấy cái này khác.}{Because it requires thinking. Not every subject does --- which is why this one feels different to you.}

\item
\studentpair{Nếu em thi lại có dễ hơn lần trước không thầy?}{If I retake the exam, will it be easier than before?}
\teacherpair{Đề không dễ hơn, nhưng nếu em học thì kết quả sẽ khác.}{The paper won't be easier --- but if you study, the result will be different.}

\item
\studentpair{Thầy có thể cho em thêm điểm vì em ngồi học rất nghiêm túc không?}{Can I get extra points for sitting very seriously in class?}
\teacherpair{Thầy sẽ cho thêm điểm khi cái sự nghiêm túc đó được phản ánh trong bài làm của em.}{I'll give extra points when that seriousness is reflected in your actual work.}

\item
\studentpair{Thầy ơi, thầy có biết ai không cần học mà vẫn điểm cao không?}{Do you know anyone who gets high scores without studying?}
\teacherpair{Thầy biết rất nhiều câu chuyện huyền thoại. Nhưng nhân vật chính thường do người kể tự bịa ra.}{I know many legends. But the main character is usually invented by the person telling the story.}

\item
\studentpair{Tại sao bài thi này lại tính một phần ba điểm tốt nghiệp ạ?}{Why does this exam count for a third of the graduation grade?}
\teacherpair{Vì hệ thống giáo dục quyết định vậy, và cả hai chúng ta đều không có tiếng nói trong chuyện đó.}{Because the education system decided so --- and neither of us has a vote in that matter.}

\item
\studentpair{Thầy nghĩ em có thể đỗ đại học không?}{Do you think I can get into university?}
\teacherpair{Thầy nghĩ em có thể, nhưng bản thân em đang làm mọi thứ để chứng minh ngược lại.}{I think you can --- but you yourself are doing everything to prove otherwise.}

\item
\studentpair{Nếu em chép bài bạn mà bài bạn đúng thì thầy có chấm điểm không?}{If I copy from a classmate and their answers are right, will you grade it?}
\teacherpair{Thầy sẽ chấm điểm. Nhưng đó là điểm của bạn em, không phải của em.}{I will grade it. But those are your classmate's points, not yours.}

\item
\studentpair{Thầy ơi, em muốn được khen dù em không làm tốt.}{I want to be praised even though I didn't do well.}
\teacherpair{Được, thầy khen em đã dũng cảm thừa nhận điều mà nhiều người cố giả vờ không phải vấn đề.}{Alright --- I praise you for having the courage to admit something most people pretend isn't a problem.}

\item
\studentpair{Học nhiều có làm hại não không thầy?}{Does studying too much harm the brain?}
\teacherpair{Học nhiều thì chưa thấy. Nhưng không học nhiều thì thầy đang thấy một số hệ quả ngay trước mắt.}{Studying too much --- haven't seen that. Not studying enough? I am observing some consequences of that right now.}

\item
\studentpair{Thầy có thể chấm thoáng hơn không ạ?}{Can you grade more leniently?}
\teacherpair{Thầy chấm đúng thay vì ``thoáng'', vì ``thoáng'' không giúp em biết mình cần sửa gì.}{I grade accurately rather than `leniently,' because `leniently' doesn't help you know what to fix.}

\item
\studentpair{Em không làm được bài vì không có cảm hứng.}{I couldn't do the assignment because I had no inspiration.}
\teacherpair{Cảm hứng là thứ xuất hiện sau khi em bắt đầu, không phải trước khi em bắt đầu.}{Inspiration appears after you start, not before.}

\item
\studentpair{Làm sao để học giỏi mà không tốn nhiều công sức ạ?}{How do I study well without putting in too much effort?}
\teacherpair{Thầy đang tìm câu trả lời cho câu đó suốt ba mươi năm dạy học và chưa ra.}{I've been searching for that answer throughout thirty years of teaching and haven't found it.}

\item
\studentpair{Thầy ơi, tại sao thầy không giải thích dễ hiểu hơn?}{Why don't you explain things more clearly?}
\teacherpair{Thầy sẽ giải thích dễ hiểu hơn khi em nghe nhiều hơn.}{I'll explain more clearly when you listen more.}

\end{enumerate}

%------------------------------------------------------------
\section{Hỏi Về Cuộc Đời, Tương Lai và Những Thứ Em Chưa Sẵn Sàng Nghe (Màn 26–50)}
%------------------------------------------------------------
\begin{enumerate}[leftmargin=*,resume]

\item
\studentpair{Thầy ơi, học xong ra đời có dễ sống không?}{When I finish school, will life be easy?}
\teacherpair{Dễ hơn nhiều so với không học xong. Thật sự mà nói.}{Much easier than if you don't finish --- honestly speaking.}

\item
\studentpair{Thầy nghĩ bao nhiêu tuổi thì nên lập gia đình?}{What age do you think is good to start a family?}
\teacherpair{Thầy nghĩ là sau khi em tốt nghiệp và có việc làm --- hai điều em đang làm chậm lại bằng cách hỏi câu này.}{I think after you graduate and have a job --- two things you are currently delaying by asking this question.}

\item
\studentpair{Nếu em không học đại học thì sao ạ?}{What if I don't go to university?}
\teacherpair{Không sao, chỉ cần em có một kế hoạch thay thế cụ thể chứ không phải là ``tính sau''.}{That's fine --- as long as you have a specific alternative plan, not just `figure it out later'.}

\item
\studentpair{Thầy có hối hận khi chọn nghề giáo không?}{Do you regret choosing to be a teacher?}
\teacherpair{Không thường xuyên. Nhưng có những buổi học như hôm nay khiến thầy nghĩ lại.}{Not usually. But some sessions like today do make me reconsider.}

\item
\studentpair{Em muốn được nổi tiếng. Thầy nghĩ sao?}{I want to be famous. What do you think?}
\teacherpair{Thầy nghĩ nổi tiếng mà không có năng lực thì ngắn. Và em cần thứ dài hơn.}{I think fame without substance is short-lived. And you need something that lasts longer.}

\item
\studentpair{Tiền hay hạnh phúc quan trọng hơn ạ?}{Which is more important --- money or happiness?}
\teacherpair{Người hỏi câu này thường không có cả hai, và nên bắt đầu từ cái nào có thể nỗ lực mà có được.}{People who ask this question usually have neither, and should start with whichever can actually be earned through effort.}

\item
\studentpair{Thầy nghĩ em sinh ra để làm gì?}{What do you think I was born to do?}
\teacherpair{Chưa rõ. Nhưng chắc chắn không phải để ngồi trong lớp mà không làm bài tập.}{Not sure yet. But certainly not to sit in class without doing homework.}

\item
\studentpair{Thầy ơi, thầy từng ước mình không làm thầy không?}{Have you ever wished you weren't a teacher?}
\teacherpair{Không thường xuyên. Nhưng mỗi lần chấm bài của em thầy lại suy nghĩ về điều đó.}{Not often. But every time I grade your work, I do reflect on it.}

\item
\studentpair{Em ghét bị kiểm soát. Thầy nghĩ sao?}{I hate being controlled. What do you think?}
\teacherpair{Thầy nghĩ đó là thái độ tốt cho người trưởng thành. Nhưng để trưởng thành thì cần phải qua giai đoạn hiện tại trước.}{I think that's a healthy attitude for an adult. But becoming an adult requires getting through the current phase first.}

\item
\studentpair{Thầy có hay tức giận học sinh không?}{Do you often get angry at students?}
\teacherpair{Thầy không tức giận. Thầy chỉ thỉnh thoảng trải qua những khoảnh khắc rất thú vị về mặt cảm xúc.}{I don't get angry. I just occasionally experience very emotionally interesting moments.}

\item
\studentpair{Làm sao để thấy vui khi đi học ạ?}{How do I enjoy going to school?}
\teacherpair{Hãy thử biết bài trước khi vào lớp. Em sẽ ngạc nhiên vì cảm giác đó tốt hơn em nghĩ.}{Try knowing the material before you arrive. You'll be surprised how much better that feels than you'd expect.}

\item
\studentpair{Thầy nghĩ em có tài năng tiềm ẩn không?}{Do you think I have hidden talent?}
\teacherpair{Thầy tin là có. Nhưng nếu nó tiềm ẩn quá lâu thì nó sẽ quen với bóng tối.}{I believe you do. But if it stays hidden too long, it will get comfortable in the dark.}

\item
\studentpair{Thầy ơi, ở tuổi thầy thì cuộc sống có ý nghĩa hơn không?}{At your age, does life feel more meaningful?}
\teacherpair{Có, vì khi lớn hơn em sẽ hiểu rằng câu hỏi đó đáng ra phải được hành động ngay ở tuổi của em, không phải chờ đến tuổi thầy.}{Yes --- because when you are older, you realize that question should have been acted on at your age, not just raised.}

\item
\studentpair{Tại sao người trẻ phải nghe lời người già?}{Why do young people have to listen to older people?}
\teacherpair{Không phải lúc nào cũng vậy. Nhưng phải nghe đủ để không lặp lại những sai lầm mà người già đã học ở giá rất đắt.}{Not always. But enough to avoid repeating mistakes that older people already paid dearly to learn.}

\item
\studentpair{Em không muốn làm việc văn phòng. Thầy hiểu không?}{I don't want an office job. Do you understand?}
\teacherpair{Thầy hiểu. Và thầy mong em tìm được con đường không phải văn phòng với đủ kỹ năng để đi được trên con đường đó.}{I do. And I hope you find a non-office path with enough skills to actually walk it.}

\item
\studentpair{Thầy có sợ chết không?}{Are you afraid of death?}
\teacherpair{Thầy sợ hơn là chết trước khi chấm hết đống bài kiểm tra kia.}{More than death, I fear dying before finishing that pile of ungraded tests.}

\item
\studentpair{Thầy có nghĩ mình hiểu học sinh không?}{Do you think you understand students?}
\teacherpair{Thầy hiểu đủ để ngạc nhiên mỗi ngày. Và đó là điều làm thầy vẫn còn muốn đi dạy.}{Enough to be surprised every single day. And that is what keeps me wanting to teach.}

\item
\studentpair{Làm thầy giáo có vui không ạ?}{Is it fun being a teacher?}
\teacherpair{Rất vui, nhất là những lúc học sinh chứng minh cho thầy thấy là nỗ lực của thầy không hoàn toàn vô ích.}{Very much so --- especially when students prove that my effort isn't entirely pointless.}

\item
\studentpair{Em muốn thay đổi thế giới. Thầy thấy được không?}{I want to change the world. Do you think I can?}
\teacherpair{Được. Nhưng trước tiên hãy thay đổi cái bàn học của em cho có trật tự cái đã.}{Certainly. But first, try changing your desk into something more organized.}

\item
\studentpair{Thầy thấy học sinh bây giờ khác ngày xưa không?}{Do you think students today are different from before?}
\teacherpair{Khác. Ngày xưa không có điện thoại. Còn câu hỏi thì vẫn giống nhau.}{Yes --- back then there were no phones. The questions, however, remain the same.}

\item
\studentpair{Bạo lực học đường sinh ra từ đâu ạ?}{Where does school bullying come from?}
\teacherpair{Từ rất nhiều nguồn. Và một trong đó là khi ai đó cảm thấy không được nghe và không biết cách xử lý điều đó.}{From many sources. One of them is when someone feels unheard and does not know how to handle that feeling.}

\item
\studentpair{Thầy nghĩ em nên yêu chưa?}{Do you think I should be in a relationship yet?}
\teacherpair{Thầy nghĩ em nên học trước, vì yêu mà không biết mình đang làm gì thì dễ gây hại cho cả hai.}{I think you should study first --- loving someone without knowing what you're doing tends to hurt both people.}

\item
\studentpair{Em thích sống tự do. Thầy có ủng hộ không?}{I like living freely. Do you support that?}
\teacherpair{Ủng hộ. Nhưng tự do mà không có kỷ luật bên trong thì chỉ là sự hỗn loạn được gọi bằng cái tên đẹp hơn.}{I support it. But freedom without internal discipline is just chaos with a prettier name.}

\item
\studentpair{Thầy có thấy em thay đổi từ đầu năm đến giờ không?}{Have you noticed me changing since the beginning of the year?}
\teacherpair{Có. Em ngủ gật khéo léo hơn một chút.}{Yes. You have gotten slightly better at dozing off discreetly.}

\item
\studentpair{Thầy ơi, lời khuyên cuối cùng của thầy cho em là gì?}{What is your final piece of advice for me?}
\teacherpair{Bắt đầu. Cái gì cũng được. Ngay bây giờ. Đừng đợi cảm hứng hay thời điểm hoàn hảo, vì cả hai sẽ không tự đến.}{Start. Anything. Right now. Don't wait for inspiration or the perfect moment --- neither will come on its own.}

\end{enumerate}

\begin{baihoc}
Câu hỏi ngây thơ nhất đôi khi là câu hỏi khó trả lời nhất. Và câu trả lời cay nhất đôi khi là lời yêu thương được bọc trong giọng gắt.\\
\textit{(The most innocent question is sometimes the hardest one to answer. And the sharpest reply is sometimes an act of care wrapped in a harsh tone.)}
\end{baihoc}

\begin{ghinhoanh}
\textbf{Học thêm tiếng Anh từ những màn đối đáp này:}\\[4pt]
\begin{itemize}[leftmargin=*]
  \item \textbf{`I understand'} --- Em hiểu rồi (thường không đúng lắm)
  \item \textbf{`Inspiration appears after you start'} --- Cảm hứng đến sau khi em bắt đầu
  \item \textbf{`Freedom without discipline is chaos'} --- Tự do không có kỷ luật là hỗn loạn
  \item \textbf{`Effort without direction hits a wall'} --- Nỗ lực không có hướng là chạy vào tường
\end{itemize}
\end{ghinhoanh}
